\chapter{REQUIREMENTS AND ANALYSIS}\label{ch:requirements}
The requirements analysis for the VeerDrishti combat vision system encompasses a comprehensive
evaluation of functional and non-functional requirements, ensuring that the system meets the
operational needs of modern warfare while adhering to technical constraints. The analysis is
structured to identify key features, performance metrics, and user interactions necessary for the
successful deployment of the system in dynamic battlefield environments.

\section{Functional Requirements}\label{sec:functional}
The functional requirements define the core capabilities of the project, including:
\begin{enumerate}[label=\roman*.]
  \item \textbf{Real-Time Object Detection:} The system must be capable of detecting and classifying
    combatants, vehicles, and weapons in real time using the integrated camera feed.
  \item \textbf{Visual Highlighting:} Detected threats should be visually highlighted on the display
    interface, providing clear and immediate feedback to the user.
  \item \textbf{Open Vocabulary Recognition:} The system should be able to identify novel or
    previously unseen threats using natural language supervision and language-vision encoders as
    a future extension; the current task variants use their configured class sets.
  \item \textbf{User Interface:} The system must provide an intuitive user interface that allows
    easy interpretation of the visual information to make informed decisions without distraction.
  \item \textbf{Data Logging:} The system must provide secure local storage of detection events and
     video metadata for post-mission analysis and system auditing in a later implementation stage.
     The present prototype records development diagnostics through Logcat rather than persistent
     mission storage.
  \label{end:functional}
\end{enumerate}
\section{Non-Functional Requirements}\label{sec:nonfunctional}
The non-functional requirements address the performance, reliability, and usability aspects of the
project, including:
\begin{enumerate}[label=\roman*.]
  \item \textbf{Low Latency:} The system must operate with minimal latency to ensure timely threat
    detection and response in high-stress combat scenarios.
  \item \textbf{Robustness:} The system should maintain high detection accuracy under varying
    environmental conditions, such as low light, occlusion, and dynamic backgrounds.
  \item \textbf{Scalability:} The system should be designed to accommodate future enhancements, such
    as additional threat categories or integration with other battlefield systems.
  \item \textbf{Security:} The system must ensure the confidentiality and integrity of the data
    collected, preventing unauthorized access and ensuring compliance with military data handling
    standards.
  \item \textbf{Usability:} The system should be user-friendly, requiring minimal training for
    effective operation, and should not contribute to cognitive overload for the user.
  \label{end:nonfunctional}
\end{enumerate}
\section{Hardware Requirements}\label{sec:hardware}
\begin{enumerate}[label=\roman*.]
  \item \textbf{Development Workstation:} High-performance laptop equipped with a dedicated GPU for
    code development, dataset preprocessing, and initial model debugging.
  \item \textbf{Edge Deployment Device:} Android smartphone with a camera and sufficient compute
    and memory to execute the deployed model and render a live preview. The exact device and
    acceleration backend must be recorded for reproducible performance measurements.
  \item \textbf{Training Resources:} GPU-equipped local or cloud resources for planned model
     fine-tuning. The final deployment record should identify the training resource and completed
     configuration used for reproducible experiments.
  \label{end:hardware}
\end{enumerate}
\section{Software Requirements}\label{sec:software}
\begingroup\emergencystretch=3em
  \begin{enumerate}[label=\roman*.]
    \item \textbf{Target Platform:} Android with camera permission, CameraX capture, and a native
      inference integration. The Android version and minimum SDK must match the application build.
    \item \textbf{Application Framework:} Flutter and Dart for the user interface and YOLOView
      integration, together with the Android native components used by the camera and model runtime.
    \item \textbf{Model Pipeline:} Ultralytics YOLO26 Nano and task variants for on-device detection,
      segmentation, and pose visualization.
    \item \textbf{Development and Diagnostics:} A Flutter/Android development environment, Android
      SDK tools, and Logcat for runtime diagnostics. Record actual package versions when preparing
      the final reproducible setup.
    \item \textbf{Model Preparation:} A compatible Ultralytics training/export environment for
      custom fine-tuning. The dataset, export format, and framework versions must be recorded with
      any completed training run.
  \end{enumerate}
\endgroup
\label{end:software}

\section{Acceptance Criteria}\label{sec:acceptance}
The requirements are considered traceable when each implemented capability can be checked through
an observable system action and a recorded test condition. Functional acceptance checks include
camera initialization, frame processing, task selection, overlay rendering, threshold control, and
graceful handling of camera or model errors. Non-functional checks should record the device,
application build, model variant, input configuration, processing time, displayed frame rate, and
operating conditions. This separation distinguishes interface verification from accuracy,
robustness, power, and sustained-performance evaluation.
\label{end:acceptance}\label{end:requirements}
